求矩阵 \(A\) 无穷范数下的条件数 首先要求矩阵 \(A\) 的逆矩阵
\(A^{- 1}\)，

\[\begin{aligned}
\begin{pmatrix}
A & \vdots & I
\end{pmatrix} & = \begin{pmatrix}
1 & - 1 & - 1 & \ldots & - 1 & - 1 & \vdots & 1 & 0 & 0 & \ldots & 0 & 0 \\
0 & 1 & - 1 & \ldots & - 1 & - 1 & \vdots & 0 & 1 & 0 & \ldots & 0 & 0 \\
0 & 0 & 1 & \ldots & - 1 & - 1 & \vdots & 0 & 0 & 1 & \ldots & 0 & 0 \\
 \vdots & \vdots & \vdots & \ddots & \vdots & \vdots & \vdots & \vdots & \vdots & \vdots & \ddots & \vdots & \vdots \\
0 & 0 & 0 & \ldots & 1 & - 1 & \vdots & 0 & 0 & 0 & \ldots & 1 & 0 \\
0 & 0 & 0 & \ldots & 0 & 1 & \vdots & 0 & 0 & 0 & \ldots & 0 & 1
\end{pmatrix} \\
 & \rightarrow \begin{pmatrix}
1 & - 1 & - 1 & \ldots & - 1 & - 1 & \vdots & 1 & 0 & 0 & \ldots & 0 & 0 \\
0 & 1 & - 1 & \ldots & - 1 & - 1 & \vdots & 0 & 1 & 0 & \ldots & 0 & 0 \\
0 & 0 & 1 & \ldots & - 1 & - 1 & \vdots & 0 & 0 & 1 & \ldots & 0 & 0 \\
 \vdots & \vdots & \vdots & \ddots & \vdots & \vdots & \vdots & \vdots & \vdots & \vdots & \ddots & \vdots & \vdots \\
0 & 0 & 0 & \ldots & 1 & 0 & \vdots & 0 & 0 & 0 & \ldots & 1 & 1 \\
0 & 0 & 0 & \ldots & 0 & 1 & \vdots & 0 & 0 & 0 & \ldots & 0 & 1
\end{pmatrix} \\
 & \ldots \\
 & \rightarrow \begin{pmatrix}
1 & 0 & 0 & \ldots & 0 & 0 & \vdots & 1 & 1 & 2 & \ldots & 2^{97} & 2^{98} \\
0 & 1 & 0 & \ldots & 0 & 0 & \vdots & 0 & 1 & 1 & \ldots & 2^{96} & 2^{97} \\
0 & 0 & 1 & \ldots & 0 & 0 & \vdots & 0 & 0 & 1 & \ldots & 2^{95} & 2^{96} \\
 \vdots & \vdots & \vdots & \ddots & \vdots & \vdots & \vdots & \vdots & \vdots & \vdots & \ddots & \vdots & \vdots \\
0 & 0 & 0 & \ldots & 1 & 0 & \vdots & 0 & 0 & 0 & \ldots & 1 & 1 \\
0 & 0 & 0 & \ldots & 0 & 1 & \vdots & 0 & 0 & 0 & \ldots & 0 & 1
\end{pmatrix} = \begin{pmatrix}
I & \vdots & A^{- 1}
\end{pmatrix}
\end{aligned}\]

不难得到这两个矩阵的范数为
\[\left\| A \right\|_{\infty} = 1 + \sum_{j = 2}^{100}\left| {- 1} \right| = 100\]
\[\left\| A^{- 1} \right\|_{\infty} = 1 + \sum_{j = 2}^{100}\left| 2^{j - 2} \right| = 2^{99}\]

所以矩阵 \(A\) 的条件数为
\[\kappa(A) = \left\| A^{- 1} \right\|_{\infty} \cdot \left\| A \right\|_{\infty} = 100 \times 2^{99}\]

教材第二章课后习题 \textbf{ }11 1 \[\begin{aligned}
\begin{pmatrix}
A & \vdots & I
\end{pmatrix} = \begin{pmatrix}
375 & 374 & \vdots & 1 & 0 \\
752 & 750 & \vdots & 0 & 1
\end{pmatrix} & \rightarrow \begin{pmatrix}
1 & \frac{374}{375} & \vdots & \frac{1}{375} & 0 \\
0 & \frac{2}{375} & \vdots & - \frac{752}{375} & 1
\end{pmatrix} \rightarrow \begin{pmatrix}
1 & \frac{374}{375} & \vdots & \frac{1}{375} & 0 \\
0 & 1 & \vdots & - 376 & \frac{375}{2}
\end{pmatrix} \\
 & \rightarrow \begin{pmatrix}
1 & 0 & \vdots & 375 & - 187 \\
0 & 1 & \vdots & - 376 & \frac{375}{2}
\end{pmatrix} = \begin{pmatrix}
I & \vdots & A^{- 1}
\end{pmatrix}
\end{aligned}\]

即

\[A^{- 1} = \begin{pmatrix}
375 & - 187 \\
 - 376 & \frac{375}{2}
\end{pmatrix}\]

所以

\[{\kappa(A)}_{\infty} = \left\| A^{- 1} \right\|_{\infty} \cdot \left\| A \right\|_{\infty} = \frac{1127}{2} \times 1502 = 846377\]

2 考虑到矩阵乘法的本质是线性变换，我们这样构造 \(b\)：

求矩阵 \(A\) 的特征值与特征向量，令
\(b = \alpha v_{1}\)，\(\delta b = \beta v_{2}\)，其中 \(v_{1}\)
为绝对值较大的特征值对应的特征向量，\(v_{2}\)
为绝对值较小的特征值对应的特征向量。

这样做的考虑是，矩阵乘法对特征向量起到拉伸的作用，超偏离拉伸效果大的方向迈一小步，会造成最终结果偏离很多，若偏向另一个特征值的方向，则效果最为显著。

\[\begin{array}{r}
\det(A - \lambda I) = \left| \begin{matrix}
375 - \lambda & 374 \\
752 & 750 - \lambda
\end{matrix} \right| = \lambda^{2} - 1125\lambda + 2 = 0 \\
 \Rightarrow \lambda_{1} = \frac{1125}{2} + \frac{\sqrt{1265617}}{2},\lambda_{2} = \frac{1125}{2} - \frac{\sqrt{1265617}}{2}
\end{array}\]

对于 \(\lambda_{1}\)，

\[\left( A - \lambda_{1}I \right)v_{1} = 0 \Rightarrow v_{1} = \begin{pmatrix}
1504 \\
 - 375 + \sqrt{1265617}
\end{pmatrix}\]

对于 \(\lambda_{2}\)，

\[\left( A - \lambda_{2}I \right)v_{2} = 0 \Rightarrow v_{2} = \begin{pmatrix}
1504 \\
 - 375 - \sqrt{1265617}
\end{pmatrix}\]

处于计算简单考虑，这里直接取 \(\alpha = \beta = 1\)，则

\[\frac{\left\| {\delta b} \right\|_{\infty}}{\left\| b \right\|_{\infty}} = \frac{1504}{1504} = 1\]

因为
\(Ax = b,A(x + \delta x) = b + \delta b \Rightarrow A\delta x = \delta b\)，解得

\[\begin{array}{r}
x = \begin{pmatrix}
 - 421873 + 375 \times \sqrt{1265617} \\
423000 - 376 \times \sqrt{1265617}
\end{pmatrix} \\
\delta x = \begin{pmatrix}
 - 634918865 - 563625 \times \sqrt{1265617} \\
636615000 + 565128 \times \sqrt{1265617}
\end{pmatrix}
\end{array}\]

则

\[\frac{\left\| {\delta x} \right\|_{\infty}}{\left\| x \right\|_{\infty}} = \frac{636615000 + 565128 \times \sqrt{1265617}}{423000 - 376 \times \sqrt{1265617}} \approx 951746993.006341\]

嗯\ldots\ldots 这里的计算可能有点问题\ldots\ldots 不过无伤大雅\ldots\ldots 吧。

不难看出，对于这样病态的矩阵 \(A\)，在
\(\frac{\left\| {\delta b} \right\|_{\infty}}{\left\| b \right\|_{\infty}}\)
较小的情况下，按照特征值选取 \(b\)、\(\delta b\) 可以使
\(\frac{\left\| {\delta x} \right\|_{\infty}}{\left\| x \right\|_{\infty}}\)
很大。

3 这一题其实与第二小题几乎一致，变形可得

\[A^{- 1}b = x,A^{- 1}(b + \delta b) = x + \delta x\]

我们只要按照矩阵 \(A^{- 1}\) 的特征向量选取 \(x\) 和 \(\delta x\)
就可以满足题意。

过程不再赘述，直接给出答案。

特征值与对应的特征向量为

\[\begin{array}{r}
\lambda_{1} = \frac{1125 + \sqrt{1265617}}{4},v_{1} = \begin{pmatrix}
 - 375 - \sqrt{1265617} \\
1504
\end{pmatrix} \\
\lambda_{2} = \frac{1125 - \sqrt{1265617}}{4},v_{2} = \begin{pmatrix}
 - 375 + \sqrt{1265617} \\
1504
\end{pmatrix}
\end{array}\]

取 \(x = v_{1},\delta x = v_{2}\)，则
\[\frac{\left\| {\delta x} \right\|_{\infty}}{\left\| x \right\|_{\infty}} = \frac{1504}{1504} = 1\]

进一步解得

\[b = \begin{pmatrix}
421871 - 375 \times \sqrt{1265617} \\
846000 - 752 \times \sqrt{1265617}
\end{pmatrix},\delta b = \begin{pmatrix}
421871 + 375 \times \sqrt{1265617} \\
752 \times \sqrt{1265617} + 846000
\end{pmatrix}\]

所以
\[\frac{\left\| {\delta b} \right\|_{\infty}}{\left\| b \right\|_{\infty}} = \frac{752 \times \sqrt{1265617} + 846000}{846000 - 752 \times \sqrt{1265617}} \approx 632810.500004216\]

\textbf{ }17 首先，计算每个平方项 \(x_{i}^{2}\) 时，有：

\[\text{ fl}\left( x_{i}^{2} \right) = x_{i}^{2}\left( 1 + \delta_{i} \right),\quad|\delta_{i}| \leq u.\]

令 \(p_{i} = \text{ fl}\left( x_{i}^{2} \right)\)，然后，计算求和
\(\sum_{i = 1}^{n}p_{i}\) 时，使用浮点加法。

假设按顺序求和，令 \(t_{1} = p_{1}\)，对于 \(k = 2\) 到 \(n\)，有：

\[t_{k} = \text{ fl}\left( t_{k - 1} + p_{k} \right) = \left( t_{k - 1} + p_{k} \right)\left( 1 + \varepsilon_{k} \right),\quad|\varepsilon_{k}| \leq u.\]

最终结果 \(\text{fl}\left( x^{T}x \right) = t_{n}\).

标准误差分析给出：

\[t_{n} = \sum_{i = 1}^{n}p_{i}\left( 1 + \eta_{i} \right),\]

其中 \(|\eta_{i}| \leq \gamma_{n - 1}\)，且
\(\gamma_{k} = \frac{ku}{1 - ku}\) 对于 \(ku < 1\)。因此，

\[t_{n} = \sum_{i = 1}^{n}x_{i}^{2}\left( 1 + \delta_{i} \right)\left( 1 + \eta_{i} \right).\]

于是，

\[\text{ fl}\left( x^{T}x \right) = \sum_{i = 1}^{n}x_{i}^{2}\left( 1 + \delta_{i} \right)\left( 1 + \eta_{i} \right) = x^{T}x + \sum_{i = 1}^{n}x_{i}^{2}\left( \delta_{i} + \eta_{i} + \delta_{i}\eta_{i} \right).\]

相对误差 \(\alpha\) 满足：

\[\text{ fl}\left( x^{T}x \right) = x^{T}x(1 + \alpha),\]

其中

\[\alpha = \frac{\sum_{i = 1}^{n}x_{i}^{2}\left( \delta_{i} + \eta_{i} + \delta_{i}\eta_{i} \right)}{x^{T}x}.\]

取绝对值：

\[|\alpha| \leq \frac{\sum_{i = 1}^{n}x_{i}^{2}|\delta_{i} + \eta_{i} + \delta_{i}\eta_{i}|}{x^{T}x} \leq \frac{\sum_{i = 1}^{n}x_{i}^{2}\left( |\delta_{i}| + |\eta_{i}| + |\delta_{i}\eta_{i}| \right)}{x^{T}x}.\]

由于 \(|\delta_{i}| \leq u\) 和 \(|\eta_{i}| \leq \gamma_{n - 1}\)，且
\(|\delta_{i}\eta_{i}| \leq u\gamma_{n - 1}\)，有：

\[|\alpha| \leq u + \gamma_{n - 1} + u\gamma_{n - 1}.\]

代入 \(\gamma_{n - 1} = \frac{(n - 1)u}{1 - (n - 1)u}\)，对于小
\(u\)，有：

\[\gamma_{n - 1} \approx (n - 1)u + O\left( u^{2} \right),\]

所以：

\[|\alpha| \leq u + (n - 1)u + O\left( u^{2} \right) = nu + O\left( u^{2} \right).\]

因此，

\[\text{ fl}\left( x^{T}x \right) = x^{T}x(1 + \alpha),\quad|\alpha| \leq nu + O\left( u^{2} \right).\]

证明 由于 \(A\)
是正定矩阵，它是对称的且所有特征值均为正实数。设其特征值为：

\[\lambda_{1} \geq \lambda_{2} \geq \ldots \geq \lambda_{n} > 0\]

矩阵 \(A\) 的 2-范数定义为：

\[|A|_{2} = \max\limits_{|x|_{2} = 1}\left| {Ax} \right|_{2}\]

对于对称矩阵，有：

\[|A|_{2} = \max(|\lambda_{i}|) = \lambda_{\max}\]

由于 \(A\) 可逆，其逆矩阵 \(A^{- 1}\) 的特征值为：

\[\frac{1}{\lambda_{1}},\frac{1}{\lambda_{2}},\ldots,\frac{1}{\lambda_{n}}\]

因此：

\[\left| A^{- 1} \right|_{2} = \frac{1}{\left| \lambda_{i} \right|} = \frac{1}{\lambda_{\min}}\]

矩阵 \(A\) 的条件数定义为：

\[\kappa(A) = |A|_{2} \cdot \left| A^{- 1} \right|_{2}\]

代入上述结果：

\[\kappa(A) = \frac{\lambda_{\max}}{\lambda_{\min}}\]
